% GENERATED FILE. Edit canonical lessons, then run npm run generate:books.
% canonical-source-hash: fnv1a64:8ee4f67f8aceb748
% canonical-lessons: KA-C538-padaka, KA-C538-kridapatu, KA-C538-vijeta, KA-C538-pandyavali, KA-C538-solisu

\chapter{Winning}
\label{ch:ka-winning}
\hlchaptermodality{538}

\begin{chapteropening}
\textbf{By the end of this chapter:} \emph{I can say a medal, an athlete, a winner, a tournament and to defeat.}

Complete the last lesson of chapter 538: I can say a medal, an athlete, a winner, a tournament and to defeat.
\end{chapteropening}

\section[\texorpdfstring{\kn{ಪದಕ} (padaka)}{ಪದಕ (padaka)}]{\kn{ಪದಕ} (padaka) — a medal}

\label{lesson:KA-C538-padaka}



\begin{quote}
Before the new one: say the Kannada for a work of art, then the Kannada for to compose.
\end{quote}

\subsection*{You'll want to know: \kn{ಪದಕ}}
\textbf{\kn{ಪದಕ}} — \emph{padaka} — \textquotedblleft{}a medal\textquotedblright{}.

\begin{quote}
\kn{ಓಟದ} \kn{ಸ್ಪರ್ಧೆಯಲ್ಲಿ} \kn{ಅವಳಿಗೆ} \kn{ಚಿನ್ನದ} \kn{ಪದಕ} \kn{ಸಿಕ್ಕಿತು}.
\end{quote}

\begin{quote}
\emph{ōṭada spardheyalli avaḷige cinnada padaka sikkitu.}
\end{quote}

\begin{quote}
She won a gold medal in the race.
\end{quote}

\kn{ಚಿನ್ನ}, \kn{ಬೆಳ್ಳಿ}, \kn{ಕಂಚು}: gold, silver, bronze.

\begin{cousinweb}
\textbf{\kn{ಪದಕ}} (\emph{padaka}) is borrowed from Sanskrit \emph{padaka}, which Monier-Williams lists as \textquotedblleft{}a kind of ornament\textquotedblright{} (Monier-Williams); a medal is an ornament you earn.
\end{cousinweb}

\begin{grammarlens}[title={What you've built}]
One more word for describing what you lived and felt.
\end{grammarlens}

\subsection*{Your turn}
\begin{itemize}
\raggedright
  \item \emph{Say it:} \emph{padaka}
  \item \emph{Say it:} \emph{padaka}, once more
  \item \emph{Say it:} \emph{ōṭada spardheyalli avaḷige cinnada padaka sikkitu.} — She won a gold medal in the race.
  \item \emph{Recall:} say the Kannada for a work of art, then the Kannada for to compose, then say \emph{padaka} again
\end{itemize}

\begin{tcolorbox}[breakable,colback=teal!4,colframe=teal!35!black,title={Before you move on}]
What does \kn{ಪದಕ} mean? (\textquotedblleft{}A medal\textquotedblright{}.) Say it once more.
\end{tcolorbox}
\section[\texorpdfstring{\kn{ಕ್ರೀಡಾಪಟು} (krīḍāpaṭu)}{ಕ್ರೀಡಾಪಟು (krīḍāpaṭu)}]{\kn{ಕ್ರೀಡಾಪಟು} (krīḍāpaṭu) — an athlete}

\label{lesson:KA-C538-kridapatu}



\begin{quote}
Before the new one: say the Kannada for to compose, then the Kannada for a medal.
\end{quote}

\subsection*{You'll want to know: \kn{ಕ್ರೀಡಾಪಟು}}
\textbf{\kn{ಕ್ರೀಡಾಪಟು}} — \emph{krīḍāpaṭu} — \textquotedblleft{}an athlete\textquotedblright{}.

\begin{quote}
\kn{ಆ} \kn{ಕ್ರೀಡಾಪಟು} \kn{ಪ್ರತಿದಿನ} \kn{ಹತ್ತು} \kn{ಕಿಲೋಮೀಟರ್} \kn{ಓಡುತ್ತಾನೆ}.
\end{quote}

\begin{quote}
\emph{ā krīḍāpaṭu pratidina hattu kilōmīṭar ōḍuttāne.}
\end{quote}

\begin{quote}
That athlete runs ten kilometres every day.
\end{quote}

Plural: \kn{ಕ್ರೀಡಾಪಟುಗಳು}, \emph{krīḍāpaṭugaḷu}.

\kn{ಕ್ರೀಡೆ}, sport + \kn{ಪಟು}, a skilled one.

\begin{cousinweb}
\textbf{\kn{ಕ್ರೀಡಾಪಟು}} (\emph{krīḍāpaṭu}) is a modern learned compound built from Sanskrit \emph{krīḍā}, \textquotedblleft{}sport, play\textquotedblright{}, and \emph{paṭu}, \textquotedblleft{}clever, skilful\textquotedblright{} (Monier-Williams).
\end{cousinweb}

\begin{grammarlens}[title={What you've built}]
One more word for describing what you lived and felt.
\end{grammarlens}

\subsection*{Your turn}
\begin{itemize}
\raggedright
  \item \emph{Say it:} \emph{krīḍāpaṭu}
  \item \emph{Say it:} \emph{krīḍāpaṭu}, once more
  \item \emph{Say it:} \emph{ā krīḍāpaṭu pratidina hattu kilōmīṭar ōḍuttāne.} — That athlete runs ten kilometres every day.
  \item \emph{Recall:} say the Kannada for to compose, then the Kannada for a medal, then say \emph{krīḍāpaṭu} again
\end{itemize}

\begin{tcolorbox}[breakable,colback=teal!4,colframe=teal!35!black,title={Before you move on}]
What does \kn{ಕ್ರೀಡಾಪಟು} mean? (\textquotedblleft{}An athlete\textquotedblright{}.) Say it once more.
\end{tcolorbox}
\section[\texorpdfstring{\kn{ವಿಜೇತ} (vijēta)}{ವಿಜೇತ (vijēta)}]{\kn{ವಿಜೇತ} (vijēta) — a winner}

\label{lesson:KA-C538-vijeta}



\begin{quote}
Before the new one: say the Kannada for a medal, then the Kannada for an athlete.
\end{quote}

\subsection*{You'll want to know: \kn{ವಿಜೇತ}}
\textbf{\kn{ವಿಜೇತ}} — \emph{vijēta} — \textquotedblleft{}a winner\textquotedblright{}.

\begin{quote}
\kn{ಸ್ಪರ್ಧೆಯ} \kn{ವಿಜೇತರಿಗೆ} \kn{ಬಹುಮಾನ} \kn{ಕೊಟ್ಟರು}.
\end{quote}

\begin{quote}
\emph{spardheya vijētarige bahumāna koṭṭaru.}
\end{quote}

\begin{quote}
They gave prizes to the winners of the competition.
\end{quote}

Plural: \kn{ವಿಜೇತರು}, \emph{vijētaru}.

\kn{ವಿಜಯ}, victory.

\begin{cousinweb}
\textbf{\kn{ವಿಜೇತ}} (\emph{vijēta}) is borrowed from Sanskrit \emph{vijetṛ}, \textquotedblleft{}a vanquisher, a conqueror\textquotedblright{} (Monier-Williams).
\end{cousinweb}

\begin{grammarlens}[title={What you've built}]
One more word for describing what you lived and felt.
\end{grammarlens}

\subsection*{Your turn}
\begin{itemize}
\raggedright
  \item \emph{Say it:} \emph{vijēta}
  \item \emph{Say it:} \emph{vijēta}, once more
  \item \emph{Say it:} \emph{spardheya vijētarige bahumāna koṭṭaru.} — They gave prizes to the winners of the competition.
  \item \emph{Recall:} say the Kannada for a medal, then the Kannada for an athlete, then say \emph{vijēta} again
\end{itemize}

\begin{tcolorbox}[breakable,colback=teal!4,colframe=teal!35!black,title={Before you move on}]
What does \kn{ವಿಜೇತ} mean? (\textquotedblleft{}A winner\textquotedblright{}.) Say it once more.
\end{tcolorbox}
\section[\texorpdfstring{\kn{ಪಂದ್ಯಾವಳಿ} (pandyāvaḷi)}{ಪಂದ್ಯಾವಳಿ (pandyāvaḷi)}]{\kn{ಪಂದ್ಯಾವಳಿ} (pandyāvaḷi) — a tournament}

\label{lesson:KA-C538-pandyavali}



\begin{quote}
Before the new one: say the Kannada for an athlete, then the Kannada for a winner.
\end{quote}

\subsection*{You'll want to know: \kn{ಪಂದ್ಯಾವಳಿ}}
\textbf{\kn{ಪಂದ್ಯಾವಳಿ}} — \emph{pandyāvaḷi} — \textquotedblleft{}a tournament\textquotedblright{}.

\begin{quote}
\kn{ಶಾಲೆಗಳ} \kn{ಕ್ರಿಕೆಟ್} \kn{ಪಂದ್ಯಾವಳಿ} \kn{ಮುಂದಿನ} \kn{ವಾರ} \kn{ಶುರುವಾಗುತ್ತದೆ}.
\end{quote}

\begin{quote}
\emph{śālegaḷa krikeṭ pandyāvaḷi mundina vāra śuruvāguttade.}
\end{quote}

\begin{quote}
The schools cricket tournament starts next week.
\end{quote}

\kn{ಪಂದ್ಯ}, a match + \kn{ಆವಳಿ}, a row: a series of matches.

\begin{cousinweb}
\textbf{\kn{ಪಂದ್ಯಾವಳಿ}} (\emph{pandyāvaḷi}) is built from \kn{ಪಂದ್ಯ}, \textquotedblleft{}a match, a contest\textquotedblright{}, and Sanskrit \emph{āvali}, \textquotedblleft{}a row, a series\textquotedblright{} (Monier-Williams).
\end{cousinweb}

\begin{grammarlens}[title={What you've built}]
One more word for describing what you lived and felt.
\end{grammarlens}

\subsection*{Your turn}
\begin{itemize}
\raggedright
  \item \emph{Say it:} \emph{pandyāvaḷi}
  \item \emph{Say it:} \emph{pandyāvaḷi}, once more
  \item \emph{Say it:} \emph{śālegaḷa krikeṭ pandyāvaḷi mundina vāra śuruvāguttade.} — The schools cricket tournament starts next week.
  \item \emph{Recall:} say the Kannada for an athlete, then the Kannada for a winner, then say \emph{pandyāvaḷi} again
\end{itemize}

\begin{tcolorbox}[breakable,colback=teal!4,colframe=teal!35!black,title={Before you move on}]
What does \kn{ಪಂದ್ಯಾವಳಿ} mean? (\textquotedblleft{}A tournament\textquotedblright{}.) Say it once more.
\end{tcolorbox}
\section[\texorpdfstring{\kn{ಸೋಲಿಸು} (sōlisu)}{ಸೋಲಿಸು (sōlisu)}]{\kn{ಸೋಲಿಸು} (sōlisu) — to defeat; to beat}

\label{lesson:KA-C538-solisu}



\begin{quote}
Before the new one: say the Kannada for a winner, then the Kannada for a tournament.
\end{quote}

\subsection*{You'll want to know: \kn{ಸೋಲಿಸು}}
\textbf{\kn{ಸೋಲಿಸು}} — \emph{sōlisu} — \textquotedblleft{}to defeat; to beat\textquotedblright{}.

\begin{quote}
\kn{ನಮ್ಮ} \kn{ಶಾಲೆಯ} \kn{ತಂಡ} \kn{ನಗರದ} \kn{ತಂಡವನ್ನು} \kn{ಸೋಲಿಸಿತು}.
\end{quote}

\begin{quote}
\emph{namma śāleya taṇḍa nagarada taṇḍavannu sōlisitu.}
\end{quote}

\begin{quote}
Our school's team beat the city team.
\end{quote}

Past (I …): \kn{ಸೋಲಿಸಿದೆ}, \emph{sōliside}.

\kn{ಸೋಲು} is to lose; \kn{ಸೋಲಿಸು} is to make someone lose.

\begin{cousinweb}
\textbf{\kn{ಸೋಲಿಸು}} (\emph{sōlisu}) is native Dravidian: the causative of \emph{sōlu}, \textquotedblleft{}to be defeated, to lose\textquotedblright{} (DEDR 3558), so \textquotedblleft{}to make someone lose\textquotedblright{}.
\end{cousinweb}

\begin{grammarlens}[title={What you've built}]
One more word for describing what you lived and felt.
\end{grammarlens}

\subsection*{Your turn}
\begin{itemize}
\raggedright
  \item \emph{Say it:} \emph{sōlisu}
  \item \emph{Say it:} \emph{sōlisu}, once more
  \item \emph{Say it:} \emph{namma śāleya taṇḍa nagarada taṇḍavannu sōlisitu.} — Our school's team beat the city team.
  \item \emph{Recall:} say the Kannada for a winner, then the Kannada for a tournament, then say \emph{sōlisu} again
\end{itemize}

\begin{tcolorbox}[breakable,colback=teal!4,colframe=teal!35!black,title={Before you move on}]
What does \kn{ಸೋಲಿಸು} mean? (\textquotedblleft{}To defeat; to beat\textquotedblright{}.) Say it once more.
\end{tcolorbox}
